\section{Proof of the Three-Distance Theorem}

We present the proof due to Liang (1979), which derives the theorem directly
from properties of the Euclidean algorithm.  Throughout, $\alpha \in (0,1)$
is irrational (the rational case is elementary).

\subsection*{Setup and Key Quantities}

Write $\alpha_0 = \alpha$ and define sequences by the Euclidean algorithm:
\[
  a_k = \lfloor \alpha_{k-1}^{-1} \rfloor, \qquad
  \alpha_k = \{{\alpha_{k-1}^{-1}}\}^{-1}\cdot\alpha_{k-1} \pmod{1},
\]
which gives the standard continued fraction $\alpha = [0; a_1, a_2, \ldots]$
with partial quotients $a_k \ge 1$.  The \emph{convergents} $p_k/q_k$ satisfy
\[
  q_{-1} = 0,\quad q_0 = 1,\quad q_{k+1} = a_{k+1} q_k + q_{k-1},
\]
and the \emph{best-approximation property}:
$\|q_k \alpha\| < \|j\alpha\|$ for all $1 \le j < q_{k+1}$.

\subsection*{The Inductive Structure}

\begin{lemma}
  For $n = q_k$ (a convergent denominator), the $n$ points of $S_n(\alpha)$
  create exactly \textbf{two} distinct gap lengths:
  \[
    \ell_k = \|q_k \alpha\| \quad \text{(short gaps, $q_{k-1}$ of them)},
    \qquad
    L_k = \|q_{k-1}\alpha\| \quad \text{(long gaps, $q_k - q_{k-1}$ of them... adjusted by $a_k$)}.
  \]
  Moreover $L_k = a_k \ell_k + \ell_{k+1}$, anticipating the next level.
\end{lemma}

\begin{proof}[Proof sketch]
  The map $T\colon x \mapsto \{x + \alpha\}$ acts as a cyclic permutation of
  $S_{q_k}(\alpha)$.  Since $T^{q_k}(x) \approx x$ (with error $\|q_k\alpha\| \ll 1$),
  the $q_k$ points are nearly equally spaced.  An exact count using
  $q_k \alpha = p_k \pm \|q_k\alpha\|$ shows exactly two gap lengths.
\end{proof}

\subsection*{The General Case}

For arbitrary $n$, choose $k$ such that $q_k \le n < q_{k+1}$ and write
$n = c \cdot q_k + r$ with $0 \le r < q_k$.

\begin{proof}[Proof of Theorem]
  We argue that the $n$ gaps of $S_n(\alpha)$ have at most three distinct lengths by
  tracking how gaps split as $n$ increases from $q_k$ to $q_{k+1}$.

  \textbf{Step 1.} At $n = q_k$, there are two gap lengths $\ell_k < L_k$.

  \textbf{Step 2.} Adding the point $\{(q_k + 1)\alpha\}, \ldots$ one at a time:
  each new point $\{(n+1)\alpha\}$ lands exactly distance $\ell_k$ from its
  predecessor $\{n\alpha\}$ (measured on $\mathbb{T}$).
  This \emph{splits a long gap $L_k$ into $\ell_k$ and $L_k - \ell_k$},
  introducing the third length $L_k - \ell_k = L_{k+1}$ (the next-level short gap).
  The sum rule holds: $L_k = \ell_k + L_{k+1}$.

  \textbf{Step 3.} After $q_{k+1} - q_k$ insertions (at $n = q_{k+1}$), all long gaps
  $L_k$ have been split, leaving only two lengths again: $\ell_{k+1} = L_k - \ell_k$
  and $\ell_k$ (now the new long gap $L_{k+1} = \ell_k$).

  At every intermediate $n$, the gaps are drawn from $\{\ell_k,\, L_{k+1},\, L_k\}$,
  a set of at most three values satisfying $L_k = \ell_k + L_{k+1}$. \qed
\end{proof}

\subsection*{Why New Points Hit Long Gaps}

The key to Step 2 is the \emph{three-distance invariant}: after $q_k$ points are placed,
the long gaps $L_k$ are each of the same length and consecutive long gaps are separated
by exactly $\ell_k$.  The rotation by $\alpha$ on $\mathbb{T}$ permutes these long gaps
cyclically (it is a cyclic permutation of the whole set $S_{q_k}$), so successive new
points $\{(q_k+j)\alpha\}$ hit successive long gaps in order, never revisiting a
short gap until all long gaps are split.
